% titel: Blatt 5, A2: Elektronenkonfiguration von Ionen
% skript: Mortimer Kap. 7 (chemie: Atombau, Elektronenkonfiguration)

Die Elektronenkonfiguration gibt an, wie die Elektronen auf die Niveaus
verteilt sind. Man unterscheidet die vollständige Schreibweise (alle Niveaus)
und die Kurzschreibweise mit Edelgas-Rumpf (z.\ B. ${}[\mathrm{Ne}]$).

\begin{werkzeugkasten}[Konfiguration von Ionen]
\begin{itemize}
  \item Aufbauprinzip, Pauli-Prinzip und Hund'sche Regel gelten unverändert.
  \item Für ein Kation werden zuerst die Elektronen der \textit{äussersten
        Schale} (grösstes $n$) entfernt --- nicht die zuletzt aufgefüllten.
        Bei Übergangsmetallen werden also erst die $ns$-Elektronen abgegeben,
        danach bei Bedarf $(n-1)d$.
  \item Für ein Anion werden so viele Elektronen addiert, dass die höchste
        Schale zur Edelgaskonfiguration ergänzt wird.
  \item Hund: Orbitaler gleicher Energie werden erst einfach besetzt. Bei
        $d^4$ und $d^9$ wird wegen der Halb-/Vollbesetzungs-Stabilität oft ein
        $s$-Elektron ins $d$ angehoben (Ausnahme beim Neutralatom).
\end{itemize}
\textit{Skript:} Mortimer Kap. 7 (Elektronenkonfiguration).
\end{werkzeugkasten}

\begin{achtung}[Darauf achten]
\begin{itemize}
  \item Abgabe-Reihenfolge nach Hauptquantenzahl, nicht nach
        Besetzungsreihenfolge: $\mathrm{[Ar]3d^6\,4s^2}$ (Fe) verliert zuerst die
        $4s$-Elektronen, nicht die $3d$.
  \item Vollständige Schreibweise muss alle Elektronen enthalten; zur Kontrolle
        zählen: Summe aller $e^-$ = Ordnungszahl minus Ladung.
  \item Voll besetzte Unterschalen ($d^{10}$) sind besonders stabil --- z.B.
        $\mathrm{Cu^+}$ mit $3d^{10}$.
\end{itemize}
\end{achtung}
